\section{Introduction}
\label{sec:introduction}

Compact objects, black holes and neutron stars, probe strong-field modifications of general relativity, among them scalar-tensor-Gauss--Bonnet gravity~\citep{TODO_StrongFieldTests}. In such theories a scalar field couples to the Gauss--Bonnet invariant $\RGB$, which in extreme curvature acts as an effective mass term for it, and the object can scalarize spontaneously: through matter for neutron stars, through curvature for black holes~\citep{Doneva2017_BH,Doneva2017_NS,DamourEspositoFarese1993}.

The field equations are non-linear, and each solution is a shooting problem with boundary conditions at $r \to \infty$ that takes seconds to minutes~\citep{TODO_SolverCost}. Bayesian parameter estimation needs many such evaluations, which makes it impractically slow~\citep{Reed2024}. Previous surrogates~\citep{Liodis2024,Kaushikk2026,Liu2026} give no procedure for building one and no domain of validity: they do not say where a prediction can be trusted.

This work gives both. Its contributions are:
\begin{enumerate}
  \item a processing pipeline as the algorithm for building surrogates of compact-object maps, which aims at the relative error the dataset allows, its data ceiling (Section~\ref{sec:pipeline});
  \item a domain of validity set by the physics, the existence edge of black holes and the stability edge of neutron stars: the surrogate never extrapolates beyond stable solutions, and refuses an input outside the domain and returns the nearest valid one (Section~\ref{sec:validity});
  \item the demonstration of the same pipeline on a two-parameter black-hole map and on a three-parameter neutron-star map, with the error of each target beside its ceiling (Section~\ref{sec:results}).
\end{enumerate}
The data science behind the design, the exploratory analysis and the three testable hypotheses $H_1$, $H_2$ and $H_3$, is in the appendix (Section~\ref{sec:hypotheses}).

Section~\ref{sec:physical-framework} sets the physical framework and its domain edges; Section~\ref{sec:pipeline} gives the pipeline; Section~\ref{sec:validity} the domain of validity; Section~\ref{sec:results} the results on both maps; Section~\ref{sec:discussion} the physical implications; Section~\ref{sec:conclusion} concludes and releases the repository.
